\section{Circuit Complexity --- Circuit Eval Run}
%%% generated by the blueprint pipeline from TCSlib.Complexity.CircuitComplexity.CircuitEvalRun
\subsection*{Circuit Complexity --- Circuit Eval Run: Introduction}

Declarations of the Lean module \texttt{TCSlib.Complexity.CircuitComplexity.CircuitEvalRun}.

\subsection*{Circuit Complexity --- Circuit Eval Run: Results}

\begin{definition}[Boolean scan update]
\lean{BoolCircuit.CircuitEval.scanF}
\label{BoolCircuit.CircuitEval.scanF}
\leanok
\uses{BoolCircuit.CircuitEval.Ph}
Define a map on a phase and a natural counter, with a Boolean input. In phase $n0$,
input true increments the counter and leaves the phase at $n0$, while input false
changes the phase to $sep$ and leaves the counter unchanged; in every other phase, the
pair is unchanged for either Boolean input.
\end{definition}

\begin{definition}[Scanning configuration]
\lean{BoolCircuit.CircuitEval.cfS}
\label{BoolCircuit.CircuitEval.cfS}
\leanok
\uses{BoolCircuit.CircuitEval.Internal.cf, BoolCircuit.CircuitEval.Internal.ip, BoolCircuit.CircuitEval.Ph, BoolCircuit.CircuitEval.St, Turing.Cfg, Turing.FinTM.bufferTape}
Let $z$ be a Boolean input list, let $\phi$ be a phase, and let $r,i$ be natural numbers
with $i\le |z|$. The associated two-work-tape machine configuration has input $z$,
first-read state carrying $\phi$, input-head position determined by $i$, an empty value
tape, a mark tape containing $r$ consecutive true cells, work-head position $r$, and
empty output.
\end{definition}

\begin{theorem}[Scanning configuration independence from proof data]
\lean{BoolCircuit.CircuitEval.cfS_congr}
\label{BoolCircuit.CircuitEval.cfS_congr}
\leanok
\uses{BoolCircuit.CircuitEval.Ph, BoolCircuit.CircuitEval.cfS}
\difficulty{1}
Let $z$ be a Boolean word, $\phi$ a scanning phase, and $r$ a natural number. If $i,i'$
are natural numbers with $i\le |z|$ and $i'\le |z|$, and $i=i'$, then the corresponding
scanning configurations are equal regardless of which bound witness is used.
\end{theorem}

\begin{theorem}[Extending a replicated tape]
\lean{BoolCircuit.CircuitEval.wr_replicate}
\label{BoolCircuit.CircuitEval.wr_replicate}
\leanok
\uses{BoolCircuit.CircuitEval.Internal.wr, Turing.FinTM.bufferTape, Turing.FinTM.bufferTape_append}
\difficulty{3}
For every natural number $r$, updating the tape that stores $r$ consecutive true symbols
starting at cell zero by setting cell $r$ to true yields the tape that stores $r+1$
consecutive true symbols starting at cell zero.
\end{theorem}

\begin{theorem}[Scanning a doubled bit string]
\lean{BoolCircuit.CircuitEval.setup_scan}
\label{BoolCircuit.CircuitEval.setup_scan}
\leanok
\uses{BoolCircuit.CircuitEval.Internal.wr, BoolCircuit.CircuitEval.LA.go, BoolCircuit.CircuitEval.Ph, BoolCircuit.CircuitEval.cfS, BoolCircuit.CircuitEval.cfS_congr, BoolCircuit.CircuitEval.evalTM, BoolCircuit.CircuitEval.lact, BoolCircuit.CircuitEval.lstep, BoolCircuit.CircuitEval.pairSym, BoolCircuit.CircuitEval.scanF, BoolCircuit.CircuitEval.wr_replicate, Turing.FinTM.bufferTape, Turing.MultiTapeTM.runFrom, Turing.MultiTapeTM.runFrom_add, Turing.MultiTapeTM.runFrom_zero, Turing.dbl}
\difficulty{6}
Let $z$ and $w$ be finite bit strings, let $r$ be a natural number, and let $\phi$ be
either the initial scan phase or the separator phase. Suppose $z$ is the concatenation
of a prefix, the string obtained by repeating each bit of $w$ twice, and a suffix.
Starting with the input head at the end of the prefix, the first work tape empty, the
second work tape containing $r$ marks from cell zero, both work heads at cell $r$, and
empty output, after $2|w|$ steps the machine is in the configuration prescribed by
folding the scan update over $w$: its phase and mark count are the resulting pair, its
input head has advanced by $2|w|$ positions, and its work tapes, heads, and output are
as initially specified.
\end{theorem}

\begin{theorem}[Scanning an initial run of ones]
\lean{BoolCircuit.CircuitEval.scanF_n0}
\label{BoolCircuit.CircuitEval.scanF_n0}
\leanok
\uses{BoolCircuit.CircuitEval.leadOnes, BoolCircuit.CircuitEval.scanF}
\difficulty{5}
Let $w$ be a Boolean word and $r\in\mathbb N$. Starting in phase $\mathrm{n0}$ with
counter $r$ and scanning $w$ from left to right, the scan ends in phase $\mathrm{sep}$
with counter $r+L$ if $w$ contains a false bit, where $L$ is the length of its initial
run of true bits; if $w$ contains no false bit, it ends in phase $\mathrm{n0}$ with
counter $r+|w|$.
\end{theorem}

\begin{definition}[Boolean interval tape]
\lean{BoolCircuit.CircuitEval.onesFrom}
\label{BoolCircuit.CircuitEval.onesFrom}
\leanok
For natural numbers $h,n$, define a tape indexed by the integers that contains true
exactly at positions $c$ satisfying $h\le c<n$, and is blank at every other position.
\end{definition}

\begin{definition}[Copying configuration]
\lean{BoolCircuit.CircuitEval.cfC}
\label{BoolCircuit.CircuitEval.cfC}
\leanok
\uses{BoolCircuit.CircuitEval.Internal.cf, BoolCircuit.CircuitEval.Internal.ip, BoolCircuit.CircuitEval.St, BoolCircuit.CircuitEval.onesFrom, Turing.Cfg, Turing.FinTM.bufferTape}
Let $z$ and $x$ be Boolean lists, let $i_0,n,h$ be natural numbers, and assume $i_0+h\le
|z|$. The associated two-work-tape machine configuration has input $z$, copy state,
input-head position determined by $i_0+h$, value tape initialized from the first $h$
entries of $x$, mark tape true at integer positions in the interval $[h,n)$ and blank
elsewhere, work-head position $h$, and empty output.
\end{definition}

\begin{theorem}[Copy configuration independence from proof data]
\lean{BoolCircuit.CircuitEval.cfC_congr}
\label{BoolCircuit.CircuitEval.cfC_congr}
\leanok
\uses{BoolCircuit.CircuitEval.cfC}
\difficulty{1}
Let $z$ be a Boolean word, $i_0,n$ natural numbers, and $x$ a Boolean word. If $h,h'$
are natural numbers with $i_0+h\le |z|$ and $i_0+h'\le |z|$, and $h=h'$, then the
copy-phase configurations constructed from these data using either bound witness are
equal.
\end{theorem}

\begin{theorem}[Copying a segment of the input]
\lean{BoolCircuit.CircuitEval.copy_run}
\label{BoolCircuit.CircuitEval.copy_run}
\leanok
\uses{BoolCircuit.CircuitEval.Internal.ip, BoolCircuit.CircuitEval.Internal.tp, BoolCircuit.CircuitEval.Internal.wr, BoolCircuit.CircuitEval.cfC, BoolCircuit.CircuitEval.cfC_congr, BoolCircuit.CircuitEval.evalTM, BoolCircuit.CircuitEval.onesFrom, BoolCircuit.CircuitEval.step_cf, BoolCircuit.CircuitEval.tr, Turing.FinTM.bufferTape_append, Turing.FinTM.inputSymbol_at, Turing.MultiTapeTM.runFrom, Turing.MultiTapeTM.runFrom_succ_eq_step, Turing.MultiTapeTM.step, Turing.dbl, Turing.length_dbl, Turing.moveInputPos_pos_of_ne_right}
\difficulty{6}
Let $e$ be an exactness flag, and let $z,\mathrm{code},x$ be Boolean words with
$z=\mathrm{dbl}(\mathrm{code})\mathbin{++}[0,1]\mathbin{++}x$. For every
$n,r,h\in\mathbb{N}$ with $h+r\le n$ and $h+r\le |x|$, and every valid bounds witnessing
$2|\mathrm{code}|+2+h\le |z|$ and $2|\mathrm{code}|+2+(h+r)\le |z|$, running the machine
for $r$ steps from the copy configuration positioned at offset $2|\mathrm{code}|+2+h$,
with the value tape holding the first $h$ bits of $x$, yields the copy configuration at
offset $2|\mathrm{code}|+2+(h+r)$, with the value tape holding the first $h+r$ bits of
$x$.
\end{theorem}

\begin{theorem}[Initial machine configuration]
\lean{BoolCircuit.CircuitEval.init_eq}
\label{BoolCircuit.CircuitEval.init_eq}
\leanok
\uses{BoolCircuit.CircuitEval.Internal.cf, BoolCircuit.CircuitEval.Internal.tp, BoolCircuit.CircuitEval.cfS, BoolCircuit.CircuitEval.evalTM, Turing.Cfg.init, Turing.FinTM.bufferTape_nil, Turing.MultiTapeTM.initCfg}
\difficulty{3}
For either Boolean mode and every Boolean input word $z$, the evaluator's initial
configuration on $z$ is the configuration in state $	ext{rd1}(	ext{n0})$ with the input
head at position $1$, both work tapes blank with both heads at position $0$, and an
empty output word.
\end{theorem}

\begin{theorem}[Rejection halts with a false output]
\lean{BoolCircuit.CircuitEval.step_rej}
\label{BoolCircuit.CircuitEval.step_rej}
\leanok
\uses{BoolCircuit.CircuitEval.Internal.cf, BoolCircuit.CircuitEval.Internal.tp, BoolCircuit.CircuitEval.LA.rej, BoolCircuit.CircuitEval.LA.toAct, BoolCircuit.CircuitEval.St, BoolCircuit.CircuitEval.evalTM, BoolCircuit.CircuitEval.step_cf, BoolCircuit.CircuitEval.tr, Turing.Cfg.inputSymbol, Turing.MultiTapeTM.runFrom, Turing.MultiTapeTM.runFrom_succ_eq_step, Turing.MultiTapeTM.runFrom_zero}
\difficulty{3}
Let $z$ be a finite bit string, and start in any state $q$ with any valid input-head
position, arbitrary two tape contents, both work heads at an arbitrary integer position
$j$, and empty output. If the machine's transition on the scanned input symbol and the
symbols under the work heads is the rejection action with input-head movement $i$, then
after one step the machine is halted and its output is the one-bit string consisting of
false.
\end{theorem}

\begin{theorem}[Reading the first input tape]
\lean{BoolCircuit.CircuitEval.step_rd1}
\label{BoolCircuit.CircuitEval.step_rd1}
\leanok
\uses{BoolCircuit.CircuitEval.Internal.cf, BoolCircuit.CircuitEval.Internal.ip, BoolCircuit.CircuitEval.Internal.wr, BoolCircuit.CircuitEval.LA.go, BoolCircuit.CircuitEval.Ph, BoolCircuit.CircuitEval.evalTM, BoolCircuit.CircuitEval.step_cf, Turing.FinTM.inputSymbol_at, Turing.MultiTapeTM.step, Turing.moveInputPos_pos_of_ne_right}
\difficulty{4}
Let $z$ be a finite bit string, let $q$ be a scan phase, and let $i$ be a natural number
such that $i+1\le |z|$. If the bit of $z$ at index $i$ is $b$, then from the
configuration in state $q$ with the input head at position $i+1$, arbitrary two
work-tape contents, both work heads at an arbitrary integer position $j$, and empty
output, one machine step moves the input head to position $i+2$ and enters the state
that records phase $q$ and the read bit $b$. The work tapes, work-head positions, and
output are unchanged.
\end{theorem}

\begin{theorem}[A bounded rejecting computation]
\lean{BoolCircuit.CircuitEval.bad_run}
\label{BoolCircuit.CircuitEval.bad_run}
\leanok
\uses{BoolCircuit.CircuitEval.Internal.ip, BoolCircuit.CircuitEval.cfS, BoolCircuit.CircuitEval.evalTM, BoolCircuit.CircuitEval.init_eq, BoolCircuit.CircuitEval.scanF, BoolCircuit.CircuitEval.setup_scan, BoolCircuit.CircuitEval.step_rd1, BoolCircuit.CircuitEval.step_rej, Turing.FinTM.bufferTape, Turing.FinTM.inputSymbol_at, Turing.MultiTapeTM.initCfg, Turing.MultiTapeTM.runFrom, Turing.MultiTapeTM.runFrom_add, Turing.MultiTapeTM.runFrom_succ_eq_step, Turing.dbl}
\difficulty{5}
Let $e$ be either Boolean value, and let $z,w,t$ be Boolean words with
$z=\mathrm{dbl}(w)\mathbin{++}t$. Suppose either $t$ is empty, $t$ consists of a single
bit, or $t$ begins with the two-bit block $10$. Then the circuit-evaluation machine with
exactness flag $e$, started on input $z$, reaches within $2|w|+2$ steps a halted
configuration whose output is the one-bit word $0$.
\end{theorem}

\begin{theorem}[Initial interval of marked cells]
\lean{BoolCircuit.CircuitEval.onesFrom_zero}
\label{BoolCircuit.CircuitEval.onesFrom_zero}
\leanok
\uses{BoolCircuit.CircuitEval.onesFrom, Turing.FinTM.bufferTape}
\difficulty{4}
For every $n\in\mathbb N$, the integer-indexed tape marking precisely the cells in
$[0,n)$ agrees at every integer cell with the tape storing a word of $n$ true bits
consecutively from cell zero and blank elsewhere.
\end{theorem}

\begin{theorem}[Empty marking at the endpoint]
\lean{BoolCircuit.CircuitEval.onesFrom_self}
\label{BoolCircuit.CircuitEval.onesFrom_self}
\leanok
\uses{BoolCircuit.CircuitEval.Internal.tapeM, BoolCircuit.CircuitEval.onesFrom}
\difficulty{3}
For every $n\in\mathbb N$, the integer-indexed tape that marks precisely the cells in
the half-open interval $[n,n)$ is the tape with no marked cells.
\end{theorem}

\begin{theorem}[A bounded halting run produces the verdict]
\lean{BoolCircuit.CircuitEval.good_run}
\label{BoolCircuit.CircuitEval.good_run}
\leanok
\uses{BoolCircuit.CircuitEval.Internal.cf, BoolCircuit.CircuitEval.Internal.cfA, BoolCircuit.CircuitEval.Internal.ip, BoolCircuit.CircuitEval.Internal.tp, BoolCircuit.CircuitEval.Internal.wr, BoolCircuit.CircuitEval.LA.rej, BoolCircuit.CircuitEval.absRun, BoolCircuit.CircuitEval.cfC, BoolCircuit.CircuitEval.cfS, BoolCircuit.CircuitEval.copy_run, BoolCircuit.CircuitEval.evalTM, BoolCircuit.CircuitEval.init_eq, BoolCircuit.CircuitEval.lact, BoolCircuit.CircuitEval.leadOnes, BoolCircuit.CircuitEval.lstep, BoolCircuit.CircuitEval.main_run, BoolCircuit.CircuitEval.onesFrom, BoolCircuit.CircuitEval.onesFrom_self, BoolCircuit.CircuitEval.onesFrom_zero, BoolCircuit.CircuitEval.rew, BoolCircuit.CircuitEval.scanF, BoolCircuit.CircuitEval.scanF_n0, BoolCircuit.CircuitEval.setup_scan, BoolCircuit.CircuitEval.step_cf, BoolCircuit.CircuitEval.step_rej, BoolCircuit.CircuitEval.tr, BoolCircuit.CircuitEval.verdict, Turing.Cfg.inputSymbol, Turing.FinTM.bufferTape, Turing.FinTM.inputSymbol_at, Turing.FinTM.timed_rewind, Turing.MultiTapeTM.initCfg, Turing.MultiTapeTM.runFrom, Turing.MultiTapeTM.runFrom_add, Turing.MultiTapeTM.runFrom_succ_eq_step, Turing.MultiTapeTM.runFrom_zero, Turing.MultiTapeTM.step, Turing.dbl, Turing.length_dbl, Turing.moveInputPos_zero, Turing.pairDecode_pairEncode, Turing.pairEncode_eq_dbl}
\difficulty{7}
For either Boolean mode, and for Boolean words $c$, $x$, and $z$ satisfying
$z=	ext{dbl}(c)\,	ext{++}\,[0,1]$ followed by $x$, there is a time $t\le 12(|z|+1)^2$ at
which the associated machine started on $z$ has halted and has output the singleton word
containing its verdict on $z$.
\end{theorem}

\begin{theorem}[Circuit evaluator runs within a quadratic bound]
\lean{BoolCircuit.CircuitEval.evalTM_computes}
\label{BoolCircuit.CircuitEval.evalTM_computes}
\leanok
\uses{BoolCircuit.CircuitEval.bad_run, BoolCircuit.CircuitEval.evalTM, BoolCircuit.CircuitEval.good_run, BoolCircuit.CircuitEval.verdict, Turing.FinTM.ComputesInTime, Turing.FinTM.computesInTime_iff, Turing.MultiTapeTM.initCfg, Turing.MultiTapeTM.runFrom, Turing.MultiTapeTM.runFrom_add, Turing.MultiTapeTM.runFrom_of_halt, Turing.eq_pairEncode_of_pairDecode, Turing.pairDecode, Turing.pairDecode_eq_none, Turing.pairEncode_eq_dbl}
\difficulty{4}
For every Boolean word $z$ and either choice of the mode parameter, the associated
two-work-tape Turing machine halts after exactly $12(|z|+1)^2$ steps, outputs the
one-symbol word whose bit is the machine's verdict on $z$, and uses some finite amount
of work-tape space.
\end{theorem}
